\section{Individual Choices}
\label{sec:individual-choices}

\subsection{Choice Probabilities (1a)}
With symmetric thresholds and an unbiased starting point, the model in
\textcite[slide 10]{woodford2026ddm} gives
\begin{equation}\label{eq:choice-probability}
  P_\alpha(\Delta)
  \;=\; \frac{e^{\alpha\Delta}}{e^{\alpha\Delta}+e^{-\alpha\Delta}}
  \;=\; \frac{1}{1+e^{-2\alpha\Delta}}, \qquad \alpha>0.
\end{equation}
The parameter $\alpha$ measures how strongly choices respond to the bid
difference. Since the two possible responses are exhaustive,
\[
  \Pr(L\mid\Delta)
  \;=\; 1-P_\alpha(\Delta)
  \;=\; \frac{e^{-\alpha\Delta}}{e^{\alpha\Delta}+e^{-\alpha\Delta}}
  \;=\; P_\alpha(-\Delta).
\]
Equal values imply a one-half probability of either response, and
reversing the value difference reverses the choice probabilities.

\subsection{Likelihood (1b)}
Let $y_i=1$ if I choose the right good and $y_i=0$ otherwise, and define
$\sigma_i=2y_i-1$. Conditional on the observed bid differences and independent
trial disturbances, the probability of the recorded sequence is
\[
  L(\alpha)
  \;=\; \prod_{i=1}^{435}
       P_\alpha(\Delta_i)^{y_i}
       [1-P_\alpha(\Delta_i)]^{1-y_i}
  \;=\; \prod_{i=1}^{435}P_\alpha(\sigma_i\Delta_i).
\]
Taking logs converts the product into a sum, following the likelihood
construction in \textcite[pp.\ 4--5]{souther2026mle}:
\begin{equation}\label{eq:individual-likelihood}
  LL(\alpha)
  \;=\; \sum_{i=1}^{435}\log P_\alpha(\sigma_i\Delta_i)
  \;=\; -\sum_{i=1}^{435}\log\!\left(1+e^{-2\alpha\sigma_i\Delta_i}\right).
\end{equation}

\subsection{Maximum Likelihood Estimate (1c)}
I maximize equation~\eqref{eq:individual-likelihood} using a logistic
regression of $y_i$ on $\Delta_i$ with no intercept. Its slope is $2\alpha$,
so dividing the estimated coefficient by two gives
\[
  \widehat{\alpha}=\AlphaIndividual,
  \qquad LL(\widehat{\alpha})=\LLIndividual.
\]
The score and curvature provide a check on this numerical estimate:
\[
  LL'(\alpha)=2\sum_i\Delta_i[y_i-P_\alpha(\Delta_i)],
  \qquad
  LL''(\alpha)=-4\sum_i\Delta_i^2P_\alpha(\Delta_i)
                         [1-P_\alpha(\Delta_i)].
\]
At the unrounded estimate, $|LL'(\widehat\alpha)|<10^{-6}$.
The second derivative is negative when at least one bid difference is nonzero,
making the positive optimum unique. Under this fit, a one-dollar advantage
for the right good raises its choice probability to about 69 percent.

\subsection{Binned Comparison (1d)}
Let $M=\max_i|\Delta_i|$, about \MaxDelta\ dollars. I divide $[-M,M]$
into 15 equal intervals, each of width $2M/15\approx\BinWidth$ dollars. Intervals include
their left endpoint and exclude their right endpoint, except that the last
also includes $M$. For the set of trials $B_j$ in bin $j$, with
$N_j=|B_j|$, I calculate
\[
  p_j=\frac{1}{N_j}\sum_{i\in B_j}y_i,
  \qquad
  e_j=\frac{1}{N_j}\sum_{i\in B_j}P_{\widehat\alpha}(\Delta_i).
\]
Inverting equation~\eqref{eq:choice-probability} gives the horizontal
coordinate,
\begin{equation}\label{eq:inverse-probability}
  d_j=P_{\widehat\alpha}^{-1}(e_j)
  \;=\;\frac{1}{2\widehat\alpha}
          \log\!\left(\frac{e_j}{1-e_j}\right).
\end{equation}
With this coordinate, $P_{\widehat\alpha}(d_j)=e_j$, so every predicted bin
point lies exactly on the fitted choice curve.

\begin{figure}[!ht]
  \centering
  \caption{\capttitle{Individual Choice Frequencies}}
  \label{fig:individual-choice}
  \includegraphics[width=0.94\textwidth]{fig-individual-choice.pdf}
  \floatnotes{Source: \texttt{indv-Data.csv}, UNI \texttt{ts3686}; 435 trials.
  Open circles show observed frequencies $p_j$ in 15 equal-width bins.
  Black dots show mean predicted probabilities $e_j$ at the inverse-probability
  coordinates in equation~\eqref{eq:inverse-probability}. The curve is
  $P_{\widehat\alpha}(d)$, with $\widehat\alpha=\AlphaIndividual$.}
\end{figure}

I almost always chose the higher-bid snack when the bids were far apart
(Figure~\ref{fig:individual-choice}). For nearly equal bids, I chose right in
52.9 percent of trials, close to the predicted 50.1 percent. The intermediate
bins are less regular. Near $d_j=1.24$, the right good won only 20 of 34
choices despite its higher bid, while the model puts its choice probability
at 73.3 percent. In the bin near $d_j=-1.86$, I chose right in 8.8 percent
of trials, roughly half the predicted 18.0 percent. Across all 15 bins,
the root mean squared difference between observed and predicted frequencies,
weighted by the number of choices in each bin, is
\ChoiceRMSEIndividual\ percentage points.
